\section{Dependence without Rivalry}

The principal objection is causal: if every created act rests upon God, divine causality seems to leave no explanatory work for created causes. That inference treats the two causes as competitors within a single order. Aquinas denies precisely that univocal comparison.

God gives creatures not only existence but definite natures and powers. Fire heats, parents generate, physicians heal, intellects understand, and wills choose. When God works in every agent, the agent's proper operation is not suppressed; it is grounded (ST I, q.~105, a.~5). Even conservation can employ ordered created causes while God remains the principal cause (q.~104, a.~2). The creature is not less real because its dependence is more radical. It can act in its own right precisely because the Creator causes it to be the kind of cause it is.

The point extends to freedom. God moves voluntary causes according to their voluntary mode rather than by converting willing into compulsion (ST I, q.~83, a.~1, ad 3; I--II, q.~10, a.~4). Divine causality does not occupy a portion of the will's agency and leave the remainder to the creature. The act, insofar as it has actuality, may be wholly from God as first cause and, in a different causal order, wholly from the person as a responsible secondary cause. When the act is sinful, its defect belongs to the created will rather than to God (ST I--II, q.~79, a.~2).

This distinction protects both metaphysics and empirical inquiry. A biochemical account of metabolism, a neurological account of perception, and a historical account of a decision can each be true at its own level. Divine causality is not inserted where an empirical explanation fails. It concerns the actual existence and operation of every cause that empirical inquiry succeeds in identifying.

\begin{studybox}{Consequences of non-rival causality}
\begin{itemize}[leftmargin=1.35em,itemsep=0.15em,topsep=0.15em]
\item The Creator is not merely a remote starter; creatures are presently conserved.
\item Created causes possess and exercise real powers rather than serving as occasions for divine action alone.
\item Dependence does not make creatures parts, modes, or pieces of the divine substance.
\item Even a hypothetically beginningless creature would depend for being; the priority is not chronological alone.
\item Primary causality can ground an act precisely as voluntary rather than replace it with force.
\item Dependence on God does not suspend moral agency before injustice, error, or any particular evil.
\end{itemize}
\end{studybox}

The \emph{Catechism} holds the two sides together. God upholds creatures ``at every moment'' and enables them to act (CCC 301); he also grants them the dignity of being causes and principles for one another (CCC 306--308). The same doctrine that denies creaturely aseity establishes creaturely agency.

The rivalry objection presupposes a self-grounded agent later displaced by divine action. Aquinas's account contains no such prior agent. Primary causality establishes the secondary cause with its proper powers; it does not enter after the fact to occupy territory the creature once possessed.

The denial of self-origination therefore entails no denial of secondary agency.
